\documentclass[11pt,a4paper]{article}
\usepackage[margin=2.4cm]{geometry}
\usepackage[T1]{fontenc}
\usepackage{lmodern}
\usepackage[ngerman]{babel}
\usepackage[hidelinks]{hyperref}
\usepackage{parskip}
\pagestyle{empty}

\begin{document}

Dr. Mach GmbH \& Co. KG\\
Am Brucker Feld 4\\
85567 Grafing

\begin{flushright}
17. August 2026
\end{flushright}

Sehr geehrte Damen und Herren,

ich bewerbe mich als Entwicklungsingenieur Medizintechnik, weil die Position die beiden Aspekte meiner bisherigen Arbeit verbindet: die Entwicklung leistungsfähiger Bildverarbeitung und ihre Anwendung in klinisch relevanten Systemen. Mein Hintergrund liegt im maschinellen Lernen für medizinische Bildgebung. Mich motiviert besonders, dass Fortschritte in der Bildverarbeitung hier nicht abstrakt bleiben, sondern medizinische Teams bei konkreten Aufgaben unterstützen können.

Mein Interesse an diesem Feld begann während meines Studiums in Projekten mit einem lokalen Krankenhaus. Die Arbeit mit Ärztinnen und Ärzten hat meinen Wunsch bestärkt, technische Lösungen für konkrete medizinische Fragen zu entwickeln. In meiner Bachelorarbeit zur anatomischen Landmarkenlokalisierung für orthopädische Anwendungen verbesserte ich in Zusammenarbeit mit Spezialisten der Wirbelsäulenmedizin ein bestehendes Modell um 30\%; die Arbeit wurde auf der DWG 2023 vorgestellt.

Anschließend arbeitete ich im dreijährigen BMBF-Projekt Photiomics an 3D-Segmentierung und Bildregistrierung für die Leberkrebsforschung. Bei RAYLYTIC entwickelte, validierte und deployte ich danach eine end-to-end 2D--3D-Registrierung zur Implantatlagenschätzung aus Röntgenbildern. Die Pipeline verbesserte die Akzeptanzrate um 45\% und reduzierte den Fehler um 15\%; heute wird das System in klinischen Studien eingesetzt.

Diese Erfahrung passt gut zu Dr. Machs Arbeit an visuellen Technologien für den Operationssaal. Ich möchte meine Kenntnisse in medizinischer Bildverarbeitung, Deep Learning und produktionsnaher Softwareentwicklung einbringen und zugleich die spezifischen Anforderungen von Kamerasystemen und klinischen Workflows weiter kennenlernen. Besonders interessant finde ich die Möglichkeit, gemeinsam mit Hardware-Teams und klinischen Partnern Lösungen zu entwickeln, bei denen Bildqualität und Benutzerfreundlichkeit unmittelbar zur Sicherheit im OP beitragen.

Ich freue mich über die Gelegenheit, im Gespräch zu erläutern, wie meine Erfahrung zur Entwicklung der nächsten Generation medizinischer Bildgebung bei Dr. Mach beitragen kann.

Mit freundlichen Grüßen\\
Kostiantyn Pysanyi

\end{document}
